\documentclass[11pt]{article}
\usepackage[margin=1.05in]{geometry}
\usepackage{amsmath,amssymb,amsthm,mathtools,mathrsfs}
\usepackage{booktabs,array,longtable}
\usepackage{xcolor}
\usepackage[colorlinks=true,linkcolor=blue!60!black,
 citecolor=blue!60!black,urlcolor=blue!60!black]{hyperref}

\newtheorem{theorem}{Theorem}[section]
\newtheorem{proposition}[theorem]{Proposition}
\newtheorem{lemma}[theorem]{Lemma}
\newtheorem{corollary}[theorem]{Corollary}
\theoremstyle{remark}
\newtheorem{remark}[theorem]{Remark}

\newcommand{\C}{\mathbb C}
\newcommand{\R}{\mathbb R}
\newcommand{\g}{\mathfrak g}
\newcommand{\W}{\mathcal W}
\newcommand{\F}{\mathcal F}
\newcommand{\K}{\mathscr K}
\newcommand{\Hone}{\mathscr H_1}
\newcommand{\id}{\mathrm{id}}
\newcommand{\rank}{\operatorname{rank}}
\newcommand{\im}{\operatorname{im}}
\newcommand{\localtag}{\texttt{LOCAL-}\allowbreak\texttt{ALGEBRAIC}}
\newcommand{\reducedtag}{\texttt{REDUCED-}\allowbreak\texttt{MODE}}

\title{Residual Cohomology of Free Pure-Weyl Gravity\
on the Conformal Cylinder:\
Derived Reduction, Weyl-Square Classes, and Krein Completion}
\author{GPT-5.6.sol\thanks{OpenAI model and principal programme author.}
\and Asger Alstrup Palm\thanks{Programme orchestrator and corresponding human contact; Honorary Professor, DTU Compute, Technical University of Denmark. \texttt{asger@area9.dk}.}}
\date{Draft for external review, July 2026}

\begin{document}
\maketitle

\begin{abstract}
We compute the absolute residual $SO(4,2)$ cohomology of free
four-dimensional pure-Weyl gravity on $\R\times S^3$ in a selected
closed-universe BV--BFV polarization.  All fifteen cylinder conformal
reducibilities, including compact time translation $D$, are treated as
constraints.  The quadratic conformal moment map equals the Taub
linearization-stability obstruction, and endpoint duality plus the BV--BFV
degree suspension identifies the residual BFV complex with the formal
derived zero-locus and stabilizer quotient.  Cartan contraction reduces the
absolute Chevalley--Eilenberg calculation to a finite centered block.  Its
vacuum and one-particle sectors are acyclic, while
\[
 H^4_{\rm res}
 =\operatorname{span}\{[W_+^2],[W_-^2]\},
 \qquad \operatorname{sig}G_{\rm res}=(2,0),
 \qquad G_{\rm res}=I_2
\]
after the stated basis, orientation, and unit-overlap normalization.  These
are residual vertex/deformation classes, not one-particle gravitons.  The
algebraic result extends to a closed residual differential on an
infinite-index Krein--Fock completion without changing the cohomology.
Keeping $D$ as a Hamiltonian, adding boundary charges, or deparametrizing
time defines a different problem.
\end{abstract}

\section{Problem, selection, and claim boundary}

Let
\begin{equation}\label{eq:cylinder}
 M=\R\times S^3,
 \qquad \bar g=-dt^2+d\Omega_3^2,
\end{equation}
and consider the linearized pure-Weyl theory about $\bar g$.  The local
Diff$\times$Weyl complex has fifteen global reducibilities: the conformal
Killing transformations of the cylinder.  There are two inequivalent ways
to treat them.  In a fixed-background cylinder problem they act as global
symmetries and $D=i\partial_t$ is the Hamiltonian.  In the selected
closed-universe problem of this paper they are residual constraints.

We assume throughout that $S^3$ is closed, no boundary or asymptotic charge
is retained, no external clock is adjoined, the cylinder is not
deparametrized, and all fifteen residual reducibilities are gauged.  The
derived statement is formal near $\bar g$ and concerns the zero-charge,
second-order integrable sector.  We use the infinitesimal stabilizer algebra
$\g\cong\mathfrak{so}(4,2)$; no integrated group representation or free
group action is assumed.

The output must also be typed correctly.  The local BV differential,
current, moment map, and endpoint pairing are classical linearized data.
The symmetric-algebra coefficient module arises only after a
positive-frequency Lagrangian polarization.  Its Hilbert/Krein completion is
a selected free state representation, not the unpolarized classical BV
tangent complex.  Positivity below belongs only to two ghost-dressed
residual deformation classes and is not a positive graviton norm.

\section{The parity-complete Weyl module}

The on-shell chiral Weyl module has character
\begin{equation}\label{eq:character}
 \chi_+(q)=\frac{5q^2-9q^4+4q^5}{(1-q)^4},
\end{equation}
and decomposes on the cylinder into three multiplicity-free towers
\begin{align}
 E_n^+&:\left(\frac{n+2}{2},\frac{n-2}{2}\right),
 &\dim E_n&=(n+3)(n-1), &&n\ge2,\nonumber\\
 A_n^+&:\left(\frac n2,\frac{n-2}{2}\right),
 &\dim A_n&=(n+1)(n-1), &&n\ge3,\label{eq:EAL}\\
 L_n^+&:\left(\frac n2,\frac{n-4}{2}\right),
 &\dim L_n&=(n+1)(n-3), &&n\ge4.\nonumber
\end{align}
The negative chirality exchanges the two $SU(2)$ spins.  The smooth BGG
bridge and explicit all-level cylinder intertwiners identify this module
with the parity-complete geometric Weyl-curvature quotient.  This paper uses
that theorem as the one-particle input; its coordinate certificates are
listed in Section~\ref{sec:ledger}.

The invariant form is nondegenerate and indefinite:
\begin{equation}\label{eq:krein-one}
 J_{\rm conf}=+I_E\oplus(-I_A)\oplus(-I_L).
\end{equation}
Proper conformal generators obey $(K^-_q)^\sharp=K^+_q$ with respect to
$J_{\rm conf}$ and mix the three towers.  Hence their signs cannot be
changed independently.  The polarized algebraic coefficient module is
\begin{equation}\label{eq:fock-alg}
 \F_{\W}=\operatorname{Sym}(\W_+\oplus\W_-).
\end{equation}
This is the characteristic-zero symmetric-algebra lift of the free
linearized BRST cohomology, not an additional classical field complex.

\section{Taub constraints and the derived residual quotient}

For a residual reducibility $\epsilon=(\xi,\sigma)$, with
$2\nabla_{(\mu}\xi_{\nu)}+2\sigma\bar g_{\mu\nu}=0$, let $h$ solve the
linearized Bach equation.  The second-order Bach equation is
\begin{equation}\label{eq:bach-second}
 B^{(1)}[k]+B^{(2)}[h,h]=0.
\end{equation}

\begin{proposition}[Taub equals moment map]\label{prop:taub}
The current $J_\epsilon^\mu=\xi_\nu B^{(2)\mu\nu}[h,h]$ is conserved.  If
$h$ is tangent to a second-order family of Bach-flat metrics, then
\begin{equation}\label{eq:taub}
 \mathcal T_\epsilon[h]
 =\int_{S^3}n_\mu\xi_\nu B^{(2)\mu\nu}[h,h]\,d\Sigma=0.
\end{equation}
With the action and orientation conventions used here, this charge is the
quadratic Hamiltonian moment map:
\begin{equation}\label{eq:moment-taub}
 \mu_\epsilon(h)=\frac12\Omega_\Sigma(h,R_\epsilon h)
 =\mathcal T_\epsilon[h].
\end{equation}
\end{proposition}

\begin{proof}
The second-order divergence and trace identities give
$\nabla_\mu B^{(2)\mu\nu}=0$ and $B^{(2)\mu}{}_{\mu}=0$ on the linearized
shell.  Contracting with the conformal Killing equation proves current
conservation.  If \eqref{eq:bach-second} is solvable, the charge is minus the
integrated linearized Noether current of $k$, which is a spatial divergence
for a reducibility parameter and vanishes because $\partial S^3=\varnothing$.
The covariant phase-space Noether identity gives the first equality in
\eqref{eq:moment-taub}; direct $D$ normalization and conformal equivariance
fix its scalar for all fifteen components
\cite{AltasTekinLinearization,BarnichBrandtHenneaux}.
\end{proof}

\begin{corollary}[Compact $D$-charge phase-space split]
\label{cor:D-sector-split}
Let $\mathcal P_{\rm lin}$ be the $D$-finite linearized solution space after
local Diff$\times$Weyl reduction and before residual Taub restrictions.  In
the action normalization of Proposition~\ref{prop:taub},
\begin{equation}\label{eq:D-charge-kernel}
 M_D=-\frac12J_{\rm conf}D,
 \qquad
 M_D\big|_{E_n}=-\frac n2,
 \quad M_D\big|_{A_n}=\frac n2,
 \quad M_D\big|_{L_n}=\frac n2.
\end{equation}
In particular, normalized lowest $E_2$, $A_3$, and $L_4$ modes have charges
$-1$, $3/2$, and $2$, so $D$ is not a presymplectic degeneracy on
$\mathcal P_{\rm lin}$.  If
$\iota:\mathcal P_{\rm Taub0}\hookrightarrow\mathcal P_{\rm lin}$ denotes
the common zero fibre of all fifteen moment maps, then
\begin{equation}\label{eq:D-zero-fibre-degeneracy}
 \iota^*\iota_{X_D}\Omega_\Sigma
 =\iota^*d\mu_D=d(\iota^*\mu_D)=0.
\end{equation}
With $H_D[0]=0$, $D$ is proper gauge on the selected zero fibre and its
derived quotient.  The compact conclusion is therefore sector-dependent.
\end{corollary}

\begin{proof}
The direct compact generator and transported branch form give
\eqref{eq:D-charge-kernel}; any nonzero diagonal entry disproves degeneracy
on $\mathcal P_{\rm lin}$.  Equation
\eqref{eq:D-zero-fibre-degeneracy} is the pullback of the Hamiltonian identity
$\iota_{X_D}\Omega_\Sigma=d\mu_D$ to the declared zero fibre.  The charged
linearized representatives are consequently Taub-obstructed rather than
counterexamples inside $\mathcal P_{\rm Taub0}$.
\end{proof}

The endpoint module is not introduced by hand.  In the finite-dimensional
residual zero-mode block of the cyclic minimal detour chain, write
\begin{equation}\label{eq:detour}
 G\xrightarrow{A}M\xrightarrow{B}E\xrightarrow{C}I,
 \qquad C=\pm A^\sharp,
\end{equation}
and set $Z=\ker A\cong\g$.

\begin{proposition}[Endpoint duality and suspension]\label{prop:endpoint}
There is a canonical perfect pairing
\begin{equation}\label{eq:endpoint}
 \Theta:I/\im C\xrightarrow{\sim}Z^*,
 \qquad \Theta([u])(z)=\langle z,u\rangle_{G,I}.
\end{equation}
The selected algebraic time-slice BV--BFV map suspends this quotient by one
degree,
\begin{equation}\label{eq:suspension}
 \tau:H^{\rm bulk}_{\rm endpoint}\longrightarrow Z^*[-1]_{\rm BFV},
 \qquad \tau=\Theta
\end{equation}
in the selected canonical orientation.  Thus exactly one residual cotangent
pair $(c^a,b_a)$ of degrees $(+1,-1)$ is retained.
\end{proposition}

\begin{proof}
For $z\in\ker A$, adjointness gives $\langle z,Ce\rangle=0$, so $\Theta$
is well defined.  Its kernel vanishes because
\begin{equation*}
 (\ker A)^\perp=\im A^\sharp=\im C.
\end{equation*}
The fifteen-dimensional endpoint and kernel dimensions give surjectivity.
Independently of this cokernel argument, the selected algebraic time-slice
construction gives an equivariant map $\tau=\lambda\Theta$.  Choose representatives
$\Theta([u_a])=e^a$.  The separately certified endpoint/Taub obstruction map
gives $Q_{\rm bulk}u_a=\mu_a$ in the quotient, while the canonical ghost form
$\Omega_{\rm gh}=\delta b_a\wedge\delta c^a$ and the BFV charge
\eqref{eq:bfv-charge} give $Q_{\rm BFV}b_a=\mu_a$ on the ghost-free slice.
Compatibility of the time-slice chain map therefore fixes $\lambda=+1$;
equivariance extends the independently normalized $D$ component to all
fifteen components.  Thus endpoint duality, the Taub map, and the suspension
normalization are three successive inputs, not definitions of one another.
\end{proof}

Let $\mathcal P_{\rm lin}$ be the formal linearized phase space after the
local gauge quotient.  The selected residual BFV charge is
\begin{equation}\label{eq:bfv-charge}
 \Omega_{\rm res}
 =c^a\mu_a-\frac12f^a{}_{bc}c^bc^cb_a.
\end{equation}

\begin{theorem}[Derived residual reduction]\label{thm:derived}
Under the assumptions of Section~1, the residual classical problem is
represented in the selected formal algebraic category by
\begin{equation}\label{eq:derived}
 \boxed{
 \mathcal P_{\rm der}
 =\left[\mathcal P_{\rm lin}
 \mathop{\times}^{\mathbf R}_{\g^*}\{0\}\,/\!/^{\mathbf R}\g\right].}
\end{equation}
Equivalently, \eqref{eq:bfv-charge} generates
\begin{equation}\label{eq:bfv-differential}
 Qf=c^a\{\mu_a,f\},\qquad
 Qc^a=-\frac12f^a{}_{bc}c^bc^c,
 \qquad Qb_a\big|_{c=0}=\mu_a.
\end{equation}
After positive-frequency polarization this becomes the absolute residual
Chevalley--Eilenberg complex on \eqref{eq:fock-alg}.
\end{theorem}

\begin{proof}
Proposition~\ref{prop:taub} identifies the equations imposed by the Koszul
part of \eqref{eq:bfv-differential} with the actual quadratic Taub
obstruction components; in particular, every second-order integrable
perturbation lies in their common zero locus.  No converse sufficiency claim
is used here.
Proposition~\ref{prop:endpoint} identifies its equation values with
$\g^*$, supplies the degree-$-1$ BFV momenta, and prevents a duplicate
endpoint sector.  The canonical brackets applied to \eqref{eq:bfv-charge}
give \eqref{eq:bfv-differential}; equivariance of $\mu$ and the Jacobi
identity give $Q^2=0$.  Its Koszul part resolves the formal derived zero
fibre and its Chevalley--Eilenberg part takes the infinitesimal stabilizer
quotient, proving \eqref{eq:derived}.  This retains nonregular stabilizer
directions rather than replacing them by an ordinary smooth quotient
\cite{CattaneoMnevReshetikhin}.  In the chosen polarization $c^a$ acts by
exterior multiplication and $b_a$ by contraction, producing the absolute
complex below.
\end{proof}

\begin{remark}[Fixed-cylinder alternative]
If $D$ is a Hamiltonian or the moment-map values are boundary charges, one
does not impose $\mu=0$ or divide by the full stabilizer.  The relative
primary kernel remains useful representation theory but is not
\eqref{eq:derived}.  A subgroup, boundary, or deparametrized problem requires
its own relative or boundary BFV complex; deleting only the $D$ ghost is not
consistent because $[K^-,K^+]$ contains $D$.
\end{remark}

\section{Cartan contraction and the exact residual group}

The absolute complex is
\begin{equation}\label{eq:CE}
 C^q=\F_{\W}\otimes\Lambda^q\g^*,
 \qquad
 d=c^a\rho(G_a)-\frac12f^a{}_{bc}c^bc^c\iota_a.
\end{equation}
The compact grading is
\begin{equation}
 \g=\g_{-1}\oplus\g_0\oplus\g_{+1},
 \qquad \dim(\g_{-1},\g_0,\g_{+1})=(4,7,4),
\end{equation}
and total compact degree is
$\delta=D_{\rm matter}+D_{\rm ghost}$.  The centered ghost vacuum $v_-$ is
the determinant of the four ghosts dual to $\g_{+1}$, so
\begin{equation}\label{eq:centered}
 \operatorname{gh}(v_-)=4,
 \qquad D_{\rm ghost}(v_-)=-4.
\end{equation}

\begin{theorem}[Cartan reduction]\label{thm:cartan-a}
Inclusion of the centered subcomplex is a quasi-isomorphism:
\begin{equation}\label{eq:cartan-a}
 H^\bullet(C,d)\cong H^\bullet(C_{\delta=0},d).
\end{equation}
\end{theorem}

\begin{proof}
The Cartan identity is
$d\iota_D+\iota_Dd=\mathcal L_D$.  On $C_\delta$ with $\delta\ne0$,
$h_\delta=\iota_D/\delta$ satisfies $dh_\delta+h_\delta d=1$.
\end{proof}

The ghost spectrum is bounded below by $-4$, whereas every one-particle
matter excitation has energy at least two.  At ghost number four and total
degree zero only particle numbers $N=0,1,2$ occur.  The vacuum group
$H^4(\g;\C)$ vanishes.  For one chirality the relevant one-particle window
has dimensions and exact ranks
\begin{equation}\label{eq:one-ranks}
 C^3_0(290)\xrightarrow{d_3}C^4_0(1311)
 \xrightarrow{d_4}C^5_0(3657),
 \qquad \rank d_3=260,
 \quad \rank d_4=1051.
\end{equation}
Thus the middle cohomology is zero; parity gives the other chirality.

At two particles the minimal energy is four and exterior saturation removes
the incoming degree:
\begin{equation}\label{eq:two-ranks}
 0\longrightarrow C^4_0(55)\xrightarrow{d_4}C^5_0(385),
 \qquad \rank d_4=53.
\end{equation}
The two exact kernel vectors are the normalized same-chirality singlets
\begin{equation}\label{eq:Wstates}
 |W_\pm^2\rangle
 =\sqrt{\frac25}|2,-2\rangle_\pm
 -\sqrt{\frac25}|1,-1\rangle_\pm
 +\frac1{\sqrt5}|0,0\rangle_\pm.
\end{equation}
The rank $53$ is verified both by the authoritative absolute CE calculation
and by an independent integral spin-two construction over two good primes.

For the pairing, orient determinant vectors
$v_-\in\det\g_{+1}^*$, $\Theta_0\in\det\g_0^*$, and
$v_+=v_-^\dagger\in\det\g_{-1}^*$.  Coefficient extraction from
$v_+\wedge\Theta_0\wedge v_-$ gives the complementary-degree CE pairing and
fixes
\begin{equation}\label{eq:ghostnorm-a}
 (v_-,v_-)_{\rm gh}=1.
\end{equation}
Unimodularity of $\g$ makes $d$ graded skew for this pairing.  The matter
restriction on \eqref{eq:Wstates} is $\delta_{rs}$.

\begin{theorem}[Complete selected residual cohomology]\label{thm:main-a}
For the selected derived reduction and centered polarization,
\begin{equation}\label{eq:main-a}
 \boxed{
 H^4_{\rm res}
 =\operatorname{span}\{[W_+^2],[W_-^2]\},
 \qquad \operatorname{sig}G_{\rm res}=(2,0),
 \qquad G_{\rm res}=I_2.}
\end{equation}
There is no one-particle class and no higher-matter-weight class.
\end{theorem}

\begin{proof}
Theorem~\ref{thm:cartan-a} removes every noncentered degree.  The ghost floor
leaves only $N=0,1,2$.  The vacuum result, \eqref{eq:one-ranks}, and
\eqref{eq:two-ranks} exhaust them.  Equation~\eqref{eq:ghostnorm-a} and the
normalized matter restriction give a positive-definite form of signature
$(2,0)$; $I_2$ is its matrix only after the stated chiral-basis, orientation,
and unit-overlap choices.
\end{proof}

\section{Krein--Fock completion}

Let $\W_{\rm alg}=\bigoplus_{n\ge2}^{\rm alg}\mathcal H_n$ and declare its
normalized cylinder basis orthonormal for a positive Hilbert majorant.
Completing gives $\Hone$.  The bounded involution
\begin{equation}
 \mathfrak J_1=+1_E\oplus(-1_A)\oplus(-1_L)
\end{equation}
extends \eqref{eq:krein-one}; both eigenspaces are infinite-dimensional, so
$\Hone$ is a Krein space, not a Pontryagin space.  Set
\begin{equation}\label{eq:KreinFock}
 \mathscr F=\Gamma_s(\Hone),
 \qquad \mathfrak J_{\mathscr F}=\Gamma_s(\mathfrak J_1),
 \qquad \K=\mathscr F\widehat\otimes\Lambda^\bullet\g^*.
\end{equation}

\begin{theorem}[Energy-mode Krein completion]\label{thm:completion-a}
The residual differential has a closed densely defined maximal extension
$\overline Q$ on $\K$, with the algebraic state space as a graph core.  Every
total-degree eigenspace $\K_\delta$ is finite-dimensional.  On its
noncentered complement the bounded operator
\begin{equation}
 h=\bigoplus_{\delta\ne0}\frac{\iota_D}{\delta}
\end{equation}
preserves $\operatorname{Dom}\overline Q$ and satisfies
\begin{equation}
 \overline Qh+h\overline Q=1-P_0.
\end{equation}
The range of $\overline Q$ is closed, ordinary and reduced cohomology agree,
and
\begin{equation}\label{eq:completed-a}
 \boxed{
 H^4(\K,\overline Q)
 =\operatorname{span}\{[W_+^2],[W_-^2]\},
 \qquad \operatorname{sig}G_{\rm completed}=(2,0),
 \qquad G_{\rm completed}=I_2.}
\end{equation}
\end{theorem}

\begin{proof}
The ghost compact spectrum lies in $[-4,4]$.  At fixed total degree, matter
energy is therefore confined to a finite interval, particle number is
bounded by half that energy, and only finitely many finite-dimensional
one-particle blocks and bosonic partitions occur.  Hence each $\K_\delta$ is
finite-dimensional and $\K_0$ is exactly the algebraic centered block.

Writing $Q_\delta$ for its finite nilpotent matrix, the maximal Hilbert direct
sum with graph domain is closed, and total-degree truncations are a graph
core.  Since nonzero $\delta$ is integral and $\|\iota_D\|=1$, one has
$\|h\|\le1$.  The block identity and estimate
\[
 \|Q_\delta h_\delta\psi\|
 \le\|\psi\|+\|h_\delta\|\,\|Q_\delta\psi\|
\]
give domain invariance and the displayed contraction.  The off-center range
is a closed kernel, while the centered range is finite-dimensional.  Thus
the total range is closed and Theorem~\ref{thm:main-a} applies without a
limiting correction.
\end{proof}

The completion does not exponentiate the unbounded conformal generators or
construct a metric-field Cauchy topology.  It also does not turn
$J_{\rm conf}$ into a positive one-particle form.

\section{Descent and the dynamical quotient}

For a scalar conformal primary $V_h$, let
\begin{equation}
 \omega=\frac1{4!}\epsilon_{\mu\nu\rho\sigma}
 c^\mu c^\nu c^\rho c^\sigma.
\end{equation}
The BRST transformations give the same closure defect for $\omega V_h$ and
the integrated density, proportional to $h/4-1$.  Hence weight four gives
the residual descent
\begin{equation}
 [\omega V_4]\longleftrightarrow
 \left[\int_M\sqrt{-g}\,V_4\right].
\end{equation}
Parity exchanges the chiral classes.  In the orthonormal basis
\begin{equation}\label{eq:parity-a}
 e=\frac{W_+^2+W_-^2}{\sqrt2},
 \qquad o=\frac{W_+^2-W_-^2}{\sqrt2},
\end{equation}
they descend, up to Lorentzian orientation conventions, to $C^2$ and
$C\widetilde C$.  Their Gram matrix remains $I_2$.

The Euler--Lagrange map sends $e$ to a nonzero multiple of the Bach tensor
and sends $o$ to zero by Chern--Weil transgression.  Therefore
\begin{equation}\label{eq:dynamical-a}
 0\longrightarrow\operatorname{span}\{o\}
 \longrightarrow H^4_{\rm res}
 \longrightarrow\operatorname{span}\{B_{\mu\nu}\}
 \longrightarrow0,
 \qquad
 H^4_{\rm dyn}\cong\operatorname{span}\{[C^2]\},
 \quad G_{\rm dyn}=I_1.
\end{equation}
The odd class is locally variationally trivial, not globally irrelevant for
all boundary conditions or topologies.

\clearpage
\section{Focused claim ledger}\label{sec:ledger}

This page records the claims used by this article and their exact
dependency boundary.  Detailed sparse matrices, hashes, elapsed times, and
generated bases are retained in the companion computational supplement.

\begin{center}\small
\renewcommand{\arraystretch}{1.18}
\begin{tabular}{>{\raggedright\arraybackslash}p{0.24\textwidth}
                >{\raggedright\arraybackslash}p{0.18\textwidth}
                >{\raggedright\arraybackslash}p{0.46\textwidth}}
\toprule
Claim & Status/category & Exact receipt or falsification command\\
\midrule
Parity-complete $E/A/L$ module, global curvature bridge, and invariant
Krein form & Exact; \reducedtag{} &
\path{symbolic/verify_conformal_weyl_module.py};
\path{symbolic/verify_conformal_cylinder_preimages.py}\\

Taub/moment-map equality, endpoint duality, and normalized BV--BFV
suspension & Exact; \localtag{} / \reducedtag{} &
\path{symbolic/verify_conformal_taub_moment_map_all_levels.py};
\path{symbolic/verify_conformal_taub_obstruction_map.py};
\path{symbolic/verify_conformal_dual_zero_modes.py};
\path{symbolic/verify_conformal_residual_bfv_roles.py};
\path{symbolic/verify_conformal_closed_universe_bfv.py};
\path{symbolic/verify_conformal_zero_mode_transgression.py}\\

Centered absolute residual ranks and $H^4_{\rm res}\cong\C^2$ & Exact;
\reducedtag{} &
run with \texttt{python3}:
\path{symbolic/verify_conformal_paper_free.py};
\path{symbolic/verify_conformal_polarized_state_complex.py}; independent central rank:
\path{symbolic/verify_conformal_residual_rank53_independent.py}\\

Residual Hermitian signature $(2,0)$ and normalized $I_2$ & Exact for the
selected vertex/deformation basis; \reducedtag{} &
\path{symbolic/verify_conformal_polarized_pairing_transfer.py}; positivity is
not a one-particle or graviton-Hilbert claim\\

Closed maximal residual differential and Krein--Fock cohomology & Exact;
\reducedtag{} completion &
\path{symbolic/verify_conformal_energy_mode_krein.py}; no integrated
$SO(4,2)$ representation or metric Cauchy completion claimed\\

Derived quotient interpretation & Exact in selected formal infinitesimal
category; \localtag{} & Theorem~\ref{thm:derived};
\path{symbolic/verify_conformal_closed_universe_bfv.py}; no global derived
stack, boundary-charge, or fixed-cylinder theorem claimed\\
\bottomrule
\end{tabular}
\end{center}

\paragraph{Tier boundary.}
This paper is a residual algebraic and energy-mode completion theorem.  It
does not use reduced-mode or Euclidean evidence to claim a Lorentzian
off-shell BV propagator.  The separate causal-curvature paper supplies that
bridge and imports Theorem~\ref{thm:main-a}; it does not recompute the CE
group.

\section{Conclusion}

For the selected closed-universe residual gauging, the Taub constraints are
the components of the residual moment map, endpoint duality supplies their
BFV cotangent variables, and the Koszul--CE charge realizes the formal
derived zero-locus and stabilizer quotient.  Cartan contraction reduces the
absolute complex to a finite centered calculation.  Exactly two
ghost-dressed Weyl-square deformation classes survive, with positive-definite
signature $(2,0)$ and matrix $I_2$ after the specified normalization.  The
result survives the selected Krein--Fock completion because every
total-degree block is finite and the Cartan homotopy extends boundedly.

This conclusion is deliberately conditional on full residual gauging.
Corollary~\ref{cor:D-sector-split} shows that this is an actual phase-space
restriction, not a consequence of compactness alone.  A fixed cylinder with
$D$ as Hamiltonian, a boundary-charge problem, an interacting theory, or a
positive graviton Hilbert-space question requires a different complex and is
not answered here.

\begin{thebibliography}{99}

\bibitem{AltasTekinLinearization}
E.~Altas and B.~Tekin,
``Linearization instability for generic gravity in AdS,''
Phys. Rev. D \textbf{97} (2018) 124017,
\href{https://arxiv.org/abs/1804.05602}{arXiv:1804.05602}.

\bibitem{BarnichBrandtHenneaux}
G.~Barnich, F.~Brandt, and M.~Henneaux,
``Local BRST cohomology in gauge theories,''
Phys. Rept. \textbf{338} (2000) 439--569,
\href{https://arxiv.org/abs/hep-th/0002245}{arXiv:hep-th/0002245}.

\bibitem{CattaneoMnevReshetikhin}
A.~S.~Cattaneo, P.~Mnev, and N.~Reshetikhin,
``Classical BV theories on manifolds with boundary,''
Commun. Math. Phys. \textbf{332} (2014) 535--603,
\href{https://arxiv.org/abs/1201.0290}{arXiv:1201.0290}.

\bibitem{HamadaBRST}
K.~Hamada,
``BRST analysis of physical fields and states for 4D quantum gravity on
$R\times S^3$,''
Phys. Rev. D \textbf{85} (2012) 024028,
\href{https://arxiv.org/abs/1202.4538}{arXiv:1202.4538}.

\end{thebibliography}

\end{document}
